\chapter{第五行星}

导航室的警告不断重复。

“第五行星稳定性丧失。地球窗口关闭倒计时启动。”

第五行星。

洛安对这个称呼毫无反应，因为对他而言，法厄同本来就是第五行星。它在火星之外，木星之内，承接内外太阳系航道，是出生、服役、战争和逃亡的全部背景。

可对我来说，这四个字像一块突然从记忆里翻出的旧石头。

太阳系没有第五行星。

至少我的时代没有。

水星、金星、地球、火星，然后是小行星带，再往外才是木星。小时候我在学校里背过，也在《另一个世界》的太空拓展包里看过。火星和木星之间只有数不清的小天体，谷神星、灶神星、智神星，还有一堆名字我记不住的石块。

可现在，星图上那里有一颗完整行星。

法厄同。

我忽然有种非常荒谬的感觉：眼前的远古战争，正同时翻开一段被我的时代删掉的天文学。

星图中央，法厄同投影开始自动展开结构图。外层穹场破裂，内层轨道环脱离，地壳裂缝像发光血管一样扩散。巨鲸坠落点被标成一个黑蓝色凹陷，凹陷周围的能量曲线正在向星球核心穿透。

“它在打穿核心。”鳞裔族女人说。

日辉族幸存者立刻反驳：“巨鲸不可能主动毁掉第五行星。它们是见证者，不是毁灭者。”

渊民族老人低声说：“见证者也会在无人肯停手时合上书。”

瑟恩看向他：“你早知道？”

老人没有回答。

导航室地面突然倾斜。所有人向一侧滑去。洛安撞上控制台，掌心巨鲸种子狠狠磕了一下。暗蓝色核亮起，一道低频沿着控制台传入星图。

星图变了。

原本的法厄同结构图外，又叠出一层更大的时间投影。

我看见无数碎片从完整行星上剥离。大陆断开，海盆蒸发，轨道环碎成闪烁的弧。法厄同在几次缓慢脉冲后彻底崩裂，变成一条围绕太阳的残骸带。

那画面和我记忆里的小行星带重叠。

不是相似。

是同一个位置。

同一个结局。

查理的现代意识几乎要从洛安身体里弹出去。我想起人类后来会说，小行星带也许从未形成完整大行星，是木星引力扰动让那些物质无法聚合。那是科学课本里的说法，冷静、干净、没有尸体，也没有巨鲸。

可这里有另一种记忆。

不是否定科学，而像是在科学能看见的结果背后，补上一段被埋掉的原因。

也许人类看到的是残骸的稳定状态。

而我现在看到的是残骸诞生时的惨叫。

“它会变成碎带。”洛安低声说。

这句话一半来自他，一半来自我。

瑟恩看了他一眼：“什么碎带？”

洛安说不出来。

因为他不该知道。法厄同在他的时代还没有变成小行星带。他只应该知道它正在死去，不应该知道它死后的名字。

日辉族幸存者忽然看向洛安。

“你身上的后来频率，来自地球之后？”

房间再次安静。

洛安握紧短刃。

瑟恩挡在他前面：“现在不是审问他的时候。”

“如果他来自未来，他知道谁继承地球。”日辉族说。

这句话像火星落进干草。

鳞裔族女人抬头。渊民族老人也看向洛安。连节肢族个体的胸部节点都开始快速闪烁。

谁继承地球？

我知道答案。

或者我以为我知道。

人类。

可这个答案忽然变得不再简单。因为站在这间导航室里的每一个族群，都像在我身体里留下过某种影子。

就在这时，巨鲸种子再次发出低频。

星图中的地球亮了起来。
